\section{Introduction}\label{sec:intro}

Whole-nervous-system connectomes invite dynamical summaries. Random walks on the wiring diagram
identify neurons that attract or repel flow \cite{lin2024}; eigen-decompositions of linear models
built from synapse counts propose ``eigencircuits'' as candidate units of computation
\cite{pospisil2024}; spiking models built from synapse counts reproduce sensorimotor responses
\cite{shiu2024}; input-normalised influence models, maximum-flow analyses and layer
assignments describe how sensory signals reach motor neurons \cite{schlegel2021,winding2023,bates2026,berg2026}.
These summaries are usually reported as point estimates computed once, without a prediction stated
in advance, without a guarantee on their numerical accuracy, and without a test of how much they
depend on the uncertainty of the reconstruction itself, although automated synapse detection
assigns every synapse a confidence score \cite{buhmann2021,berg2026}.

The complete connectome of the male \Dros{} central nervous system \cite{berg2026,malecnsdata},
which joins the brain and the ventral nerve cord (VNC), makes a whole-CNS version of these questions
possible. (A female brain-and-cord connectome, BANC, appeared shortly before \cite{bates2026}; it was
analysed with an unsigned influence model, spectral clustering and centrality, not with the mixing
questions asked here.) We ask a simple question of the synapse-flow random walk, in which a walker
at a neuron moves to a postsynaptic partner with probability proportional to the number of
synapses: after how many synaptic steps has the walk forgotten where it started? The quantity that
answers it is the Dobrushin ergodicity coefficient $\Dob(P^r)$ of the $r$-step chain
\cite{dobrushin1956,seneta2006}, the largest total-variation distance that two starting neurons can
still keep after $r$ steps \cite{levin2017}. The number of steps after which it becomes small is the
classical counterpart of the entanglement-breaking index of a quantum channel, the number of
self-compositions after which a channel forgets its input \cite{kopel2026a}. Unlike a spectral gap, $\Dob(P^r)$ at finite $r$ can be bounded from both sides
by certified numerical bounds, and the bounds can be carried through a stated class of
reconstruction errors.

We did this as a pre-registered study \cite{nosek2018}. An exploratory analysis of the male CNS
(19 September 2026) produced five hypotheses with numerical pass and fail criteria, about the
spectral gap of the chain and its dependence on the neck connective (H1), the localisation of the
leading modes of an input-normalised signed linear map (H2), the certified mixing depth at a
threshold of five synapses per connection (H3), the robustness of the certified depth of the
one-synapse chain to per-synapse confidence (H4), and the dependence of the signed spectrum on the neurotransmitter sign convention
(H5). The pre-registration was made public on 26 September 2026 \cite{kopel2026reg}, the
confirmatory scripts were tested on synthetic data only and frozen \cite{kopel2026freeze}, and each
hypothesis was then run once. Section~\ref{sec:confirm} reports all five against their frozen
criteria: H2 and H5 held, H1, H3 and H4 failed. Section~\ref{sec:explore} reports analyses that we
made afterwards to understand the failures; they are exploratory and labelled as such. They show
that the slowest mode of the chain at the five-synapse threshold is carried by eleven neurons of the
cord, seven of them efferent neurons of the mesothoracic neuromere, that become a near-absorbing set
once weak connections are dropped, that the same kind of boundary set carries the slowest mode of the FlyWire female brain at one
synapse \cite{dorkenwald2024,schlegel2024}, and that chains built from the confidence data change the depth
profile far less than the worst-case certificate allows, most of whose slack lies in the step-by-step
accumulation of perturbations rather than in the radius assigned to each neuron.
